% ============================================================
% Paper Draft Content (Template-Agnostic)
% Usage:
%   - Copy sections into ACL/EMNLP template later.
%   - Keep labels/tables; replace placeholders with final numbers.
% ============================================================

\section{Introduction}
\label{sec:intro}
Legal question answering in Vietnamese remains constrained by limited high-quality, domain-grounded datasets, especially for multimodal educational materials (text + figures/tables). We present a practical data construction pipeline for Vietnamese legal QA from textbook-like sources and evaluate whether quality-controlled synthesis yields measurable gains over simpler generation settings.

This work targets legal education use cases and focuses on a reproducible pipeline that transforms raw legal learning documents into structured, evaluable QA records. Our central claim is that multimodal, quality-controlled synthesis can improve dataset quality and downstream legal QA utility compared with text-only or no-filter alternatives.

\paragraph{Contributions.}
\begin{itemize}
\item A Vietnamese legal multimodal QA data construction pipeline with staged quality controls.
\item A controlled experimental matrix with baselines and ablations tailored to compute-constrained settings.
\item A human evaluation protocol with inter-annotator agreement reporting.
\item A reproducible artifact package (splits, integrity checks, quality gates, run commands).
\end{itemize}

\section{Related Work}
\label{sec:related}
We situate this work at the intersection of legal NLP benchmarks, synthetic data generation, and multimodal/document QA. Prior legal benchmarks (e.g., LexGLUE, LegalBench, COLIEE-style tasks) demonstrate broad legal reasoning challenges but provide limited coverage for Vietnamese multimodal educational legal QA. Synthetic data approaches improve scale but require careful quality validation. LLM-as-judge frameworks are practical yet known to exhibit bias, motivating cross-judge checks and human auditing. Document-grounded QA literature further motivates explicit handling of OCR and visual grounding.

\section{Task Definition and Research Questions}
\label{sec:task_rq}
\paragraph{Task.}
Given legal educational source documents, construct a QA dataset with fields:
\texttt{qa\_id, domain\_tag, bloom\_level, context\_payload, question\_content, is\_multimodal, candidate\_answers, ground\_truth, legal\_rationale}.

\paragraph{Research Questions.}
\begin{itemize}
\item \textbf{RQ1:} Does the proposed pipeline produce higher-quality QA data than text-only and no-filter baselines?
\item \textbf{RQ2:} Which components contribute most to quality (multimodal grounding, chunk cleaning, legal-rationale prompting, filtering thresholds)?
\item \textbf{RQ3:} Does dataset-driven adaptation improve downstream legal QA performance on a held-out human gold test set?
\end{itemize}

\paragraph{Hypotheses.}
\begin{itemize}
\item \textbf{H1:} Multimodal grounding improves groundedness and reduces unsupported questions.
\item \textbf{H2:} Quality filtering improves legal fluency and reduces factual errors.
\item \textbf{H3:} Using the constructed dataset improves downstream legal QA metrics over non-adapted baselines.
\end{itemize}

\section{Dataset Construction Pipeline}
\label{sec:pipeline}
\subsection{Stage Overview}
The pipeline includes: (1) PDF digitization to markdown + images, (2) context structuring and optional chunk cleaning, (3) QA generation with Bloom-level control and legal rationale constraints, and (4) quality evaluation/filtering plus final JSONL export.

\subsection{Data Splits and Leakage Control}
We use document-level split assignment (train/dev/test) to prevent chunk leakage across splits. Split artifacts and chunk manifests are generated with fixed seeds and stored for reproducibility.

\subsection{Quality Gates}
Before experimental runs, we enforce four gates:
\begin{itemize}
\item Gate A: OCR/chunk coverage
\item Gate B: valid context ratio
\item Gate C: parse-valid QA ratio
\item Gate D: duplicate QA rate
\end{itemize}

\section{Experimental Setup}
\label{sec:setup}
\subsection{Systems Compared}
\paragraph{Primary system.}
Full multimodal pipeline with quality filtering.

\paragraph{Baselines.}
\begin{itemize}
\item \textbf{B1} Text-only generation (no visual context in generation prompt).
\item \textbf{B2} No-judge/no-filter setting (raw generation retained).
\item \textbf{B3} Cross-judge subset evaluation with a different judge model family.
\end{itemize}

\paragraph{Ablations.}
\begin{itemize}
\item \textbf{A1} Disable multimodal alignment criterion during filtering.
\item \textbf{A2} Disable chunk cleaning in structuring stage.
\item \textbf{A3} Remove explicit legal-syllogism prompting.
\item \textbf{A4} Threshold sensitivity (strict vs softer thresholds on dev).
\end{itemize}

\subsection{Metrics}
\paragraph{Intrinsic metrics.}
Groundedness, legal fluency, multimodal alignment, parse-valid rate, duplicate QA rate, diversity.

\paragraph{Human metrics.}
Legal correctness, groundedness, question quality, distractor quality, visual alignment; agreement reported with percent agreement and Cohen's $\kappa$.

\paragraph{Downstream metrics.}
Task-specific correctness (e.g., EM/F1 or exact/partial legal QA correctness), with 95\% confidence intervals via bootstrap.

\subsection{Compute and Reproducibility}
We target limited GPU availability and prioritize reproducible, fixed-split experiments. All runs log seeds, commands, model configs, and hardware summaries.
To reduce self-judge bias, evaluation uses a judge model configuration separated from the generation model by default; same-model judging is explicitly blocked unless manually overridden.

\section{Results}
\label{sec:results}
\paragraph{Claim--Evidence policy (mandatory).}
Every added claim in this paper must include: (i) motivation/rationale, (ii) empirical evidence, and (iii) reproducible trace (artifact path, command, or table reference). Claims without supporting evidence are treated as hypotheses or future work, not conclusions.

\subsection{Dataset Statistics}
\begin{table}[t]
\centering
\small
\begin{tabular}{lrr}
\hline
Split & \#Docs & \#Contexts \\
\hline
Train & 26 & 4183 \\
Dev   & 11 & 897 \\
Test  & 11 & 900 \\
\hline
\end{tabular}
\caption{Document-level split summary from generated split artifacts (seed=42), after pruning context documents not present in raw sources.}
\label{tab:split_stats}
\end{table}

\begin{table}[t]
\centering
\small
\begin{tabular}{lc}
\hline
Quality Gate & Value \\
\hline
Gate A (OCR/chunk coverage) & 106.25\% \\
Gate B (valid context ratio) & 100.00\% \\
Gate C (parse-valid QA ratio) & 100.00\% \\
Gate D (duplicate QA rate) & 0.00\% \\
\hline
\end{tabular}
\caption{Automatic quality gates from pre-experiment checks.}
\label{tab:quality_gates}
\end{table}

\subsection{Main Baseline/Ablation Comparison}
\begin{table*}[t]
\centering
\small
\begin{tabular}{lcccccc}
\hline
System & Grounded. $\uparrow$ & Legal Fluency $\uparrow$ & MM Align. $\uparrow$ & Parse-valid $\uparrow$ & Dup. Rate $\downarrow$ & Human Score $\uparrow$ \\
\hline
Primary (current, pre-Stage4) & \texttt{N/A} & \texttt{N/A} & \texttt{N/A} & 100.00\% & 0.00\% & \texttt{TBD} \\
B1 Text-only & \texttt{TODO} & \texttt{TODO} & \texttt{TODO} & \texttt{TODO} & \texttt{TODO} & \texttt{TODO} \\
B2 No-filter & \texttt{TODO} & \texttt{TODO} & \texttt{TODO} & \texttt{TODO} & \texttt{TODO} & \texttt{TODO} \\
A1 No MM-align & \texttt{TODO} & \texttt{TODO} & \texttt{TODO} & \texttt{TODO} & \texttt{TODO} & \texttt{TODO} \\
A2 No cleaning & \texttt{TODO} & \texttt{TODO} & \texttt{TODO} & \texttt{TODO} & \texttt{TODO} & \texttt{TODO} \\
A3 No syllogism & \texttt{TODO} & \texttt{TODO} & \texttt{TODO} & \texttt{TODO} & \texttt{TODO} & \texttt{TODO} \\
\hline
\end{tabular}
\caption{Main comparison table. Current row uses available intrinsic metrics from 120 QA records in chunk checkpoints; quality and human-dependent metrics are filled after full Stage 3/4 and annotation.}
\label{tab:main_results}
\end{table*}

\paragraph{Current intrinsic snapshot.}
Using currently available QA chunk checkpoints (fallback source when \texttt{raw\_qa\_pairs.json} is empty), we observe: 120 QA records, parse-valid rate 100.00\%, duplicate rate 0.00\%, diversity ratio 88.33\%, and Bloom-level distribution of 30 (Remember), 47 (Understand), 43 (Apply). These are preliminary and will be superseded by full-run baseline/ablation results.

\subsection{Human Evaluation Agreement}
\begin{table}[t]
\centering
\small
\begin{tabular}{lcc}
\hline
Criterion & Percent Agreement & Cohen's $\kappa$ \\
\hline
Legal correctness & \texttt{TODO} & \texttt{TODO} \\
Groundedness & \texttt{TODO} & \texttt{TODO} \\
Question quality & \texttt{TODO} & \texttt{TODO} \\
Distractor quality & \texttt{TODO} & \texttt{TODO} \\
Visual alignment & \texttt{TODO} & \texttt{TODO} \\
\hline
\end{tabular}
\caption{Inter-annotator agreement on human evaluation subset.}
\label{tab:agreement}
\end{table}

\subsection{Downstream Utility}
We compare a base downstream setup against the same setup adapted with the constructed dataset. Report absolute/relative gains and 95\% confidence intervals.

\subsection{Claim-to-Evidence Ledger}
\label{sec:claim_evidence}
\begin{table*}[t]
\centering
\small
\begin{tabular}{p{0.19\linewidth}p{0.20\linewidth}p{0.19\linewidth}p{0.34\linewidth}}
\hline
Claim & Why it matters & Evidence type & Evidence pointer \\
\hline
Split prevents leakage & Fair evaluation & Integrity/leakage audit & \texttt{research/artifacts/context\_manifest.json}; \texttt{research/results/integrity\_*.json} \\
Quality-controlled pipeline improves data quality & Core contribution & Baseline/ablation + human eval & Table~\ref{tab:main_results}, Table~\ref{tab:agreement} \\
Dataset improves downstream QA & Utility contribution & Downstream metrics + CI & \texttt{research/results/downstream\_*.json}, bootstrap report \\
\hline
\end{tabular}
\caption{Claim-to-evidence ledger. Every major statement must map to auditable evidence.}
\label{tab:claim_evidence}
\end{table*}

\section{Error Analysis}
\label{sec:error}
Provide at least 30 representative failure cases, grouped by taxonomy:
\begin{itemize}
\item unsupported or weakly grounded question
\item legal reasoning inconsistency
\item distractor ambiguity
\item visual misalignment
\item OCR-induced context corruption
\end{itemize}

\section{Discussion}
\label{sec:discussion}
Discuss where improvements are robust, where gains are fragile, and how compute constraints affect external validity.

\section{Ethics, Licensing, and Responsible Release}
\label{sec:ethics}
Describe source permissions, release constraints, misuse risks, and mitigation measures. Explicitly state that dataset outputs are not legal advice.

\section{Limitations}
\label{sec:limitations}
Cover domain coverage limits, synthetic label noise risk, judge bias, and dependence on OCR quality.

\section{Conclusion}
\label{sec:conclusion}
Summarize findings against RQ1--RQ3 and outline next steps for broader legal-domain expansion and stronger human-centered validation.

% ============================================================
% Appendix Suggestions
% A. Prompts and parsing rules
% B. Split manifests and integrity checks
% C. Command list for reproducing each table
% D. Additional error cases
% ============================================================
